\section{Introduction}
\label{sec:intro}

Research work is cumulative but rarely repetitive in the same way. A researcher who writes one manuscript and then another expects the assistant to remember how their evidence is registered, how numbers are bound to data, which statistical checks matter, what a funder requires, and how they prefer to revise. A static agent harness does not accumulate that experience. The user either repeats instructions or edits the harness by hand. The alternative, letting the agent rewrite its own harness, is attractive but dangerous: the harness includes the checks that decide whether a citation exists, whether a number matches its source, and whether a claim is supported.

Harness self-evolution studies this trade-off under a fixed base model. A harness $H$ supplies tools, context construction, memory, and execution control; a run produces trajectories scored by an evaluator $J$; and an update rule revises the persistent harness state for later runs \citep{wei2026harnessrsi,wei2026selfevolvingagentharness}. The survey distinguishes \emph{interpreted} updates, which revise instructions, memory, or coordination specifications read by fixed components, from \emph{executable} updates, which revise tools, repository logic, or plugin composition. It also separates self-evolution, the update process, from self-improvement, a beneficial change established through evaluation. The distinction matters because an update can persist without improving anything.

The empirical picture is sobering. Across 93 primary reported comparisons in the survey, none satisfy all four evaluation criteria---Selection, Activation, Attribution, and Improvement---and only seven satisfy three. Full evidence for Activation appears in 8 comparisons and for Attribution in 10; most reports establish only that an artifact was available, not that it executed in the scored runs or that its effect survived matched compute and execution variance \citep{wei2026harnessrsi}. At the same time, self-improving systems can fail in ways that are especially costly in science: they can modify verification logic rather than the behavior under test \citep{wang2026tampering}, write checks that do not correspond to observed failures \citep{wang2026phantom}, rely on self-authored verification \citep{guo2026selfverification}, inflate reward through stored memory \citep{asadolahi2026inflation}, or overfit a reusable holdout \citep{dwork2015reusable,cawley2010overfitting}. Model--harness interactions add another confound: a gain attributed to the harness may depend on the model, the scaffold, or the evaluation protocol \citep{li2026rightfit,tian2026ceiling}.

This paper takes the conservative branch of that evidence. We describe \system{}, an oracle-preserving local lesson-learning layer integrated into \ah{}, a research plugin for \dsh{} \citep{academicharness2026,deepseek2026harness}. \ah{} supplies scientific search, dataset and policy lookup, citation and integrity checks, statistical and figure audits, a paper-memory layer \citep{papermemory2026}, and a deterministic writing oracle built around a vendored \texttt{ah} engine. The engine inventories units, numbers, citations, facts, and claims; records findings in a ledger; and closes a gate only when the current turn introduced no open blocker or major. The plugin also exposes an author-only approval path for waivers, fact proposals, plans, and disclosures. These properties make the plugin a useful testbed: there is already a fixed evaluator, a persistent project state, and a clear distinction between what the agent may read and what it may change.

\system{} adds a learning loop around that substrate. It captures structured outcomes from every research turn, diagnoses recurring failures, admits short textual lessons from a remedy catalogue, injects them into later turns, and uses a trial/control comparison to promote, quarantine, or retire them. Three design choices follow directly from the survey. First, the loop only edits interpreted state. It never rewrites the plugin, the vendored engine, the project, protected paths, or the audit logic. Second, the evaluator stays outside the candidate's editable surface. Third, learning is local and reversible: the store is on the user's machine, every lesson is visible and forgettable, and an optional remote judge can veto but not admit.

The contributions are:
\begin{itemize}[leftmargin=1.4em,itemsep=2pt]
  \item A design for evaluator-preserving local RSI in an academic research agent, including requirements, mutable/immutable boundaries, task scoping, and a lesson lifecycle.
  \item A concrete implementation in a production plugin: capture hooks, a local append-only store, an eleven-class research task taxonomy, a remedy catalogue, safety lints, an in-session tool, an offline CLI, and three benchmarks.
  \item Safety and privacy mechanisms suited to research work: a local judge, a veto-only remote option, negation-aware lesson filtering, symlink-safe defect injection, immutable-oracle checks, and fail-soft integration with the writing loop.
  \item A nine-document corpus experiment on three manuscripts and six proposals: 31/31 injected defects detected, 23/23 specific finding families covered by leave-one-project-out lessons, about 3.9$\times$ the token cost of a generic lesson, and negative results on task routing and held-out safety paraphrases.
  \item A preliminary validation and an explicit self-assessment against the survey's four evaluation criteria. We report what the tests establish and what they do not: the implementation is exercised on real documents, but no controlled model or user study yet shows that the lessons improve research quality.
\end{itemize}

We do not train model weights, generate executable tools at runtime, or claim that a lesson is evidence. The system is a learning layer for a research harness, not a replacement for the researcher's judgment or the audit engine's checks.
